\documentclass[english]{tudscrartcl}

\usepackage{graphicx}
\graphicspath{
	{./resources/}
	{./resources/tud_logo_svg_rgb/Blau/}
	{./resources/tud_logo_svg_rgb/Schwarz/}
	{./resources/tud_logo_svg_rgb/Weiss/}
	{./resources/plots/}
	{./resources/plots/acae}
	{./resources/plots/native}
}

% This was taken from tudscr src to change the TUD logo to the new one (TUD Corporate Design)
\makeatletter
\renewcommand*\tud@mainlogo@set{ % changed to override the exiting cmd
  \if@tud@mainlogo@set%
    \def\@tempa{\includegraphics[keepaspectratio,width=\tud@dim@logowidth]}%
    \tud@savelayerbox{main@black}{\@tempa{TUD_Logo_RGB_horizontal_schwarz_en}}%
    \tud@savelayerbox{main@HKS41}{\@tempa{TUD_Logo_RGB_horizontal_blau_en}}%
    \tud@savelayerbox{main@white}{\@tempa{TUD_Logo_RGB_horizontal_weiss_en}}%
    \settoheight\tud@dim@mainlogoheight{\tud@uselayerbox{main@black}}%
    \global\tud@dim@mainlogoheight=\tud@dim@mainlogoheight%
    \gdef\tud@mainlogo@wrn##1{%
      \ifdim##1<\ta@bcor\relax%
        \ClassWarning{\TUD@Class@Name}{%
          The selected page layout means that the\MessageBreak%
          logo of TUD extends beyond the printing area. \MessageBreak%
          The inner margin is smaller than BCOR\MessageBreak%
          (`BCOR=\the\ta@bcor', inner margin is \the##1)\MessageBreak%
          Maybe you should decrease the current value\MessageBreak%
          of DIV (`DIV=\the\ta@div')%
        }
        \global\let\tud@mainlogo@wrn\@gobble%
      \fi%
    }%
    \tud@headlogo@set%
    \global\@tud@footlogo@option@settrue%
    \tud@footlogo@option@set%
    \global\@tud@mainlogo@setfalse%
  \fi
}
\makeatother
% ~This was taken from tudscr src to change the TUD logo to the new one (TUD Corporate Design)

\usepackage{babel}
\usepackage[babel]{microtype}
\usepackage[utf8]{inputenc}
\usepackage[T1]{fontenc}
\usepackage{clrscode}
\usepackage{xcolor}
\usepackage{tabularx}
\usepackage{array}
\usepackage{xurl}
\usepackage{subcaption}
\usepackage{acronym}
\usepackage{hyperref}
\hypersetup{
    colorlinks,
    citecolor=black,
    filecolor=black,
    linkcolor=black,
    urlcolor=black
}
\usepackage[backend=biber,style=numeric,sorting=none]{biblatex}
\addbibresource{references.bib}

% smart metering
\acrodef{SG}{smart grid}
\acrodef{AMI}{advanced metering infrastructure}
\acrodef{MDMS}{meter data management system}
\acrodef{CE}{central entity}
\acrodef{SM}{smart meter}
\acrodef{SMGW}{smart meter gateway}

% hardware
\acrodef{HW}{hardware}
\acrodef{MCU}{microcontroller unit}
\acrodef{CPU}{central processing unit}
\acrodef{RAM}{random access memory}
\acrodef{NVIC}{nested vectored interrupt controller}
\acrodef{MPU}{memory protection unit}
\acrodef{FPU}{floating-point unit}
\acrodef{RTL}{register-transfer-level}
\acrodef{ADC}{analog-to-digital converter}
\acrodef{RTC}{real-time clock}
\acrodef{RMS}{root-mean-square}
\acrodef{ISA}{instruction set architecture}

\acrodef{PLL}{phase-locked loop}
\acrodef{HSE}{high speed external}
\acrodef{UART}{Universal Asynchronous Receiver-Transmitter}

% software
\acrodef{OS}{operating system}
\acrodef{RT-OS}{real-time operating system}
\acrodef{STD}{C standard library}

% simulation
\acrodef{ISS}{instruction-set-simulator}
\acrodef{TB}{translation block}

% benchmark
\acrodef{DMIPS}{Dhrystone mega instructions per second}
\acrodef{EEMBC}{embedded microprocessor benchmark consortium}

% acae
\acrodef{ACAE}{ARM cycle accounting emulator}

\begin{document}

	\addtokomafont{subject}{\sffamily}
	\title{Project title: from task desc.}
	\author{Jakob Fregin}
	\date{August 2026}

  	\maketitle
  	\pagenumbering{arabic}

	%\clearpage
	
	\section{TODO}
		- double col 6-8p. \\
		- technical report, less backroung more refs. \\
		- content: \\
			what was done?
			how was it done?
			how to repoduce?
			why even, what's notable (e.g. compared to our existing measurements)
			results: some proof of work, no you don't need to calculate hashes on paper and hope for zeros
		- goal:
			- (me) a broader dataset usable for cycle model improvements based on workloads comparable to real world deployments
			- as question: How does the performance gap of acae translate to real world representative workloads
			- plots:
				- each 
			- future work:
				- maybe we can even provide some guidelines to performance measurements to improve acae
				- explicitly which configurations represent proper measurements and which to avoid
				

		Wenn dein Ziel also z.B. ist, die Relevanz von Embench im Vergleich mit Coremark festzustellen, wären Differenzen zwischen den beiden Benchmarks spannend.
		Oder wenn ich mit Embench z.B. den zeitlichen Ablauf besser erforschen kann, dann wäre ein Plot für die Zeit gut. Usw....

	\section{STACK}
		SECTIONS: \\
		- Intro / Motivation / bit of Background / our work \\
		- Results

		Motivation: \\
		ACAE currently deviates from real HW performance as proven by (ref ACEA) using Dhrystone and CoreMark
		In order to improve ACEA to be useful for accurate performance measurements for e.g. rapid prototyping or scientific research both omitting the use of HW,
		a broader view on comparable performance data is needed.
		This data should be measured using close to real world workloads, to ensure that measured loads are represented properly by our cycle model.
		\\
		Embench!
		is a suite
		provides broader workloads
		is explicitly built for embedded systems
		built on BEEBs
		aims to provide real world workloads
		runs on hot caches to represent real world workloads better -> (WILD CLAIM - SOURCE?) e.g. SHA256 performance is comparable to real deployments
		\\
		Work: \\
		using ACAE in its state as of presented in (ref ACEA) + slight compatibility changes (WHICH? - check git log + diff)
		we integrated embench as a benchmark suite
		\\
		CHANGES:
		acae:\\
		- adjusted Make-System to accomodate benchmark suites by adding the "list" target providing a list of all benchmarks a suite offers
		embench:\\
		- discarded the python harness -> acae provides its own build system
		\\

	\section{Introduction}
		\label{sec:intro}
		Something profound goes here. \\
	
	\printbibliography

\end{document}
